\chapter{Considerações Finais}
\label{cap:consid}

Este trabalho partiu da arquitetura de \citeonline{Souza2023}, uma CNN 1D com a fusão do identificador do superpixel SEEDS, e a otimizou em quatro eixos (organização do código, paralelismo, automação de hiperparâmetros e segmentadores), medindo cada ganho contra a configuração original, com a fusão. A rede foi reimplementada em PyTorch, em um pacote modular e determinístico, e sobre essa base foram implementadas três estratégias de paralelismo, três estratégias de busca de hiperparâmetros e quatro algoritmos de segmentação, avaliados em Indian Pines, University of Pavia e Salinas Valley.

\section{Objetivos}

O primeiro objetivo, reimplementar a arquitetura de forma modular e reprodutível, foi atendido: a reimplementação preserva a arquitetura, o PCA, a divisão dos dados e os hiperparâmetros de treino, troca a regra de parada antecipada (\autoref{subsec:desTreino}) e produz sempre o mesmo resultado para a mesma configuração. Os valores reportados por \citeonline{Souza2023} ficam acima dos da reimplementação por menos de 1 ponto de OA em Indian Pines e Salinas Valley, 1,6 ponto em University of Pavia e 7,2 pontos de AA em Indian Pines, em uma execução de cada lado (\autoref{tab:linhaBase}). A organização também revelou três defeitos (\autoref{sec:resOrganizacao}): o não determinismo do LSC, o gerador aleatório compartilhado entre threads e o piso de componentes do PCA.

O segundo, explorar algoritmos de segmentação, foi atendido com SEEDS, SLIC, LSC e Watershed sob otimização conjunta (\autoref{sec:resSegmentadores}). O segmentador é uma das escolhas de maior impacto: na validação média dos 540 \textit{trials}, separa 14,89 pontos, atrás apenas do \textit{dropout} e da taxa de aprendizado (\autoref{fig:eixos}). Nenhum segmentador superou o SEEDS de forma afirmável, e a pureza das regiões não acompanhou o resultado (\autoref{tab:estruturaSeg}); o número de amostras de treino por região é maior no SEEDS, o de melhor resultado, mas não ordena os outros três, e quatro segmentadores de uma única cena não demonstram uma causa.

O terceiro, o ajuste automático de hiperparâmetros, foi atendido com busca aleatória, TPE e algoritmo genético, que substituíram a busca em grade e a Otimização Bayesiana com processo gaussiano previstas no projeto (\autoref{sec:desHPO}, \autoref{sec:resHPO}). Com o SEEDS, a busca elevou a OA de Indian Pines de 0,7441 para 0,8174, com 10\% do treino e 5 sementes, um ganho de 7,33 pontos de OA, 10,81 de AA e 8,52 de Kappa (\autoref{tab:impactoHPO} e \autoref{tab:impactoHPOCompleto}); o ganho com o treino completo não foi medido, e o efeito de cada mudança da configuração vencedora não foi isolado. Nenhuma estratégia se separou do ruído entre sementes (\autoref{tab:estrategias}), e o orçamento de épocas por \textit{trial} pesou mais que a estratégia.

O quarto, o processamento paralelo, foi atendido em três níveis (\autoref{sec:resParalelismo}). Entre os ramos do pré-processamento, o ganho chegou a 1,22$\times$, no teto do desenho, apenas em Salinas Valley (\autoref{tab:ramos}), e o pré-processamento é 0,9\% da execução. A partição da cena com processos e pool aquecido acelerou o treino em até 4,47$\times$ com 16 partições em Indian Pines (\autoref{tab:poolQuente}), enquanto as threads, limitadas pelo GIL, ficaram mais lentas que a execução sequencial a partir de 4 partições. A fusão concorrente, implementada com uma variável de condição e validada em 723 execuções com as mesmas classes do lote, atingiu o teto dado pela segmentação em University of Pavia e Salinas Valley e chegou a 2,01$\times$ em Indian Pines contra o lote com uma thread do PyTorch (\autoref{tab:fusao}); contra o lote com 16 threads, rende de 1,01$\times$ a 1,29$\times$, e com uma posição por tarefa fica de 7 a 35 vezes mais lenta.

O quinto, a análise comparativa por métricas de acurácia e custo computacional, corresponde ao \autoref{cap:resultados}, que reporta OA, AA, Kappa, tempo, memória e \textit{speedup} contra referências explícitas. Com os objetivos específicos cumpridos, foi atendido o objetivo geral de analisar o impacto de técnicas de segmentação, parametrização automática e processamento paralelo na classificação de imagens hiperespectrais com CNNs 1D.

\section{Hipóteses}

\begin{enumerate}
    \item \textit{Diferentes algoritmos de segmentação apresentarão desempenhos variados em termos de acurácia e tempo de processamento, sendo possível identificar alternativas ao SEEDS para cenários específicos.} Confirmada em parte: os desempenhos variam, mas nenhuma alternativa superou o SEEDS em acurácia de forma afirmável. Em University of Pavia, o Watershed segmentou cerca de 60 vezes mais rápido que o SEEDS (\autoref{tab:estruturaSeg}), o que o torna candidato quando a segmentação pesa no tempo total, como na inferência de uma cena nova.
    \item \textit{A otimização automática de hiperparâmetros identificará configurações superiores ao ajuste manual, melhorando métricas de acurácia (OA, AA, Kappa) de forma reprodutível.} Confirmada para a configuração de referência, com o SEEDS em Indian Pines e 10\% do treino; para os demais segmentadores, o ganho de OA ficou dentro do ruído entre sementes (\autoref{tab:impactoHPO}).
    \item \textit{A otimização conjunta de segmentação e CNN superará a otimização isolada de cada componente.} Não testada como formulada, porque a otimização isolada de cada componente não foi executada. Os resultados mostram que a comparação com a CNN fixa é inválida: com a configuração do SEEDS, o SLIC obteve Kappa de 0,03 em University of Pavia e, otimizado para si, empatou com o SEEDS.
    \item \textit{A paralelização do processamento espacial e espectral resultará em redução mensurável do tempo de execução em relação à implementação sequencial, explorando a independência de dados entre as operações e a execução concorrente via threads.} Confirmada com ressalvas: as duas estratégias com threads reduziram o tempo (até 1,22$\times$ e 2,01$\times$), mas sobre etapas pequenas de uma execução dominada pelo treino; a maior redução, de 4,47$\times$, veio da partição com processos.
\end{enumerate}

\section{Limitações}

A busca de hiperparâmetros e a comparação de segmentadores usaram 10\% do treino, e o ganho da busca com o treino completo não foi medido. A comparação de estratégias, o ganho sobre a configuração original, a partição da cena e a fusão concorrente com uma posição por tarefa foram medidos apenas em Indian Pines. A linha de base compara uma execução de cada lado, e a ablação do identificador usa 3 sementes, sem desvio registrado. As medições de tempo foram feitas em uma única máquina, cujo número de núcleos e memória condicionam os \textit{speedups}.

\section{Trabalhos Futuros}

\begin{itemize}
    \item \textbf{Fusão por vetores aprendidos}: substituir o identificador numérico por um vetor aprendido por região, que elevou a OA de 0,8981 para 0,9471 em Indian Pines (\autoref{tab:ablacao}), e avaliá-lo com mais sementes e com os quatro segmentadores;
    \item \textbf{Paralelismo sem GIL}: repetir a fusão concorrente e a partição em threads em uma versão do CPython sem a trava global \cite{PEP703};
    \item \textbf{Processos com memória compartilhada}: medir o método \textit{forkserver} e o compartilhamento de memória entre processos \cite{PyTorchMultiprocessing}, para reduzir os cerca de 481 MB por processo;
    \item \textbf{Paralelismo na busca}: executar em paralelo os \textit{trials} independentes, como os da busca aleatória ou os de uma geração do algoritmo genético;
    \item \textbf{Ampliação da campanha}: medir o ganho da busca nas outras cenas e com o treino completo, executar a otimização isolada de cada componente e separar, por ablação, os efeitos do \textit{dropout}, do lote e do PCA.
\end{itemize}
